% Cover letter — Reliability Engineering & System Safety (Elsevier)
% Submission portal: https://submit.elsevier.com/JRESS
\documentclass[11pt,a4paper]{article}
\usepackage[T1]{fontenc}
\usepackage[utf8]{inputenc}
\usepackage[a4paper,margin=2.5cm]{geometry}
\usepackage[UKenglish]{babel}
\usepackage{hyperref}
\hypersetup{hidelinks}
\setlength{\parskip}{0.6em}
\setlength{\parindent}{0pt}

\begin{document}

\noindent 25 May 2026

\bigskip

\noindent The Editor-in-Chief\\
\emph{Reliability Engineering \& System Safety}\\
Elsevier

\bigskip

\noindent Dear Editor,

We submit for your consideration the manuscript entitled ``BowTie Risk
Analysis of the Valve Train in a High-Performance VW AP~1.8 Cylinder
Head Operating at 10{,}000~rpm'' as a Research Paper.

The manuscript proposes a quantitative-coupled BowTie procedure for
design-stage reliability assessment of the valve train of a
high-performance spark-ignition cylinder head, organised around a
top event of catastrophic significance for high-rpm engines: the loss
of follower contact between valve and cam at 10{,}000~rpm. Beyond the
qualitative bow diagram, each preventive barrier is associated with a
deterministic indicator computed from the analytical sizing of the
intake valve and the dual chrome--silicon valve springs (dynamic
follower-contact ratio, natural-frequency margin against cam
harmonics, coil-bind clearance, and Goodman fatigue factor of safety),
with literature-anchored pass criteria. A Monte Carlo refinement
under ASTM~A877/A877M tolerances ($N=2\times 10^{5}$ trials) converts
the deterministic at-risk classification into probabilities of barrier
compliance failure under manufacturing-realistic dispersion, and
surfaces a third marginal barrier that the nominal-point analysis
would have missed. The procedure is qualitatively validated against
four documented valve-train and engine-component failure cases from
the recent literature, and a concrete experimental validation
roadmap --- spring-rig modal verification, motored-engine speed
sweep, and inner-spring fatigue campaign, with explicit acceptance
criteria --- is specified for empirical confirmation.

We believe \emph{Reliability Engineering \& System Safety} is the
right venue for this work for three reasons. First, the BowTie
methodology is a central thread of the journal: from the seminal
contributions of Khakzad, Khan and Amyotte (RESS~2012,
DOI~10.1016/j.ress.2012.04.003) to recent extensions by Vileiniskis
and Remenyte-Prescott (RESS~2017), Taleb-Berrouane and colleagues
(RESS~2021), Xie and colleagues (RESS~2021), and Sezer and colleagues
(RESS~2023), RESS has consistently published the methodological
advances on which the present manuscript builds. Second, the
quantitative coupling between barrier and deterministic indicator
addresses a gap explicitly identified in the BowTie literature ---
how to convert a qualitative diagram into a verifiable engineering
check at the design stage --- and is therefore aligned with the
journal's mission of advancing methods for the safety and reliability
of complex engineered systems. Third, the application domain
(valve-train dynamics at extreme rotational speed in a passenger-car
platform with a wide installed base in Brazil) extends BowTie risk
analysis from its classical chemical/process-safety home into the
mechanical-reliability domain, where the literature is currently
thinner and where the proposed deterministic-indicator coupling has
immediate engineering value.

The manuscript is original, has not been published elsewhere, and is
not under consideration by any other journal. All four authors have
approved the submission, and the authors declare no competing
financial interests or personal relationships that could have
influenced the work. No funding was received for this research. No
new experimental data were generated; the Monte Carlo script is
provided as supplementary material.

We respectfully suggest the following independent reviewers as
potentially appropriate. None has co-authored work with the present
authors, and none is affiliated with the Universidade Cat\'olica de
Petr\'opolis (UCP) or the Pontif\'icia Universidade Cat\'olica do Rio
de Janeiro (PUC-Rio). Institutional contact details can be located on
each researcher's institutional page.

\begin{enumerate}
\item \textbf{Prof.\ Valerio Cozzani} --- Department of Civil,
Chemical, Environmental and Materials Engineering (DICAM), University
of Bologna, Italy. Leading authority on safety-barrier performance,
LOPA, and Natech-cascade analysis; multiple recent RESS publications
on barrier degradation (e.g.\ RESS~2021, DOI~10.1016/j.ress.2021.107634).
Relevance: barrier-based reasoning is the central methodological
device of the present manuscript.

\item \textbf{Prof.\ Chia-Fen Chi} --- Department of Industrial
Management, National Taiwan University of Science and Technology
(NTUST), Taipei, Taiwan. Recent RESS work on FMEA classification of
passenger-vehicle recalls (RESS~2020,
DOI~10.1016/j.ress.2020.106929). Relevance: the present manuscript
explicitly compares the quantitative-coupled BowTie against
FMEA/RPN, including an AIAG--VDA 2019 scoring of the six identified
threats, which is directly within Prof.\ Chi's published expertise.

\item \textbf{Prof.\ Saeed Khodaygan} --- Department of Mechanical
Engineering, Sharif University of Technology, Tehran, Iran. Recent
RESS work on Bayesian-reliability-based multi-objective tolerance
design of mechanical assemblies (RESS~2021,
DOI~10.1016/j.ress.2021.107748). Relevance: the manuscript's Monte
Carlo refinement under ASTM~A877/A877M wire-and-strength tolerances
sits squarely in the mechanical-reliability-under-tolerance domain.

\item \textbf{Prof.\ Frank Guldenmund} --- Section Safety and Security
Science, Faculty of Technology, Policy and Management, Delft
University of Technology, Netherlands. Co-author of the canonical
BowTie review (de Ruijter and Guldenmund, \emph{Safety Science},
2016, DOI~10.1016/j.ssci.2016.03.001) cited in our manuscript;
senior authority on the qualitative formulation of the BowTie
diagram.

\item \textbf{Prof.\ Gabriele Landucci} --- Department of Civil and
Industrial Engineering, University of Pisa, Italy. Co-author of
multiple RESS papers on safety-barrier performance under
domino-effect and Natech scenarios (e.g.\ RESS~2020,
DOI~10.1016/j.ress.2019.106597). Relevance: barrier-performance
assessment under realistic dispersion is the methodological core of
our Monte Carlo layer.
\end{enumerate}

The final selection of reviewers is, of course, entirely at the
editor's discretion.

We look forward to your editorial decision and remain available for
any clarification.

\bigskip

\noindent Yours sincerely,

\bigskip\bigskip

\noindent Darlan Costa Porto\\
Corresponding author\\
Programa de Mestrado em Sistemas de Engenharia\\
Universidade Cat\'olica de Petr\'opolis\\
Petr\'opolis, RJ, Brazil\\
\href{mailto:darlan.42540004@ucp.br}{darlan.42540004@ucp.br}

\end{document}
